\documentclass[11pt]{article}
\usepackage[margin=1in]{geometry}
\usepackage{amsmath,amssymb}
\usepackage{array}
\usepackage{booktabs}
\usepackage{longtable}
\usepackage[T1]{fontenc}

\title{GUTCP Cosmological Model 2023: Audit of a Redshift Counterpart}
\author{}
\date{\today}

\newcommand{\dd}{\mathrm{d}}
\newcommand{\zG}{z_{\mathrm{GUTCP}}}
\newcommand{\RG}{\mathcal{R}}
\newcommand{\EG}{\mathcal{E}}
\newcommand{\ZRel}{\mathfrak{Z}_{\mathrm{GUTCP}}}
\newcommand{\HRel}{\mathfrak{H}_{\parallel,\mathrm{GUTCP}}}
\newcommand{\DRel}{\mathfrak{D}_{\mathrm{GUTCP}}}
\newcommand{\ARel}{\mathfrak{A}_{\mathrm{GUTCP}}}
\newcommand{\BRel}{\mathfrak{B}_{\mathrm{GUTCP}}}

\begin{document}
\maketitle

\section*{Scope}

This note audits the reconstructed 2023 Chapter 32 material in
\texttt{docs/reconstructed.tex.pages/} against the existing Mathematica notebook
\texttt{GUTCPCosmologicalModel.nb}.  The target is not a FLRW redshift
parameter.  The target is the smallest GUTCP-native candidate for a wavelength
shift relation that can be written from Mills's equations without importing
\(1+z=a(t_o)/a(t_e)\).

The notation below treats cosmological observables as relations.  A relation
may later become a singleton graph after branch, domain, and calibration
constraints are added, but that uniqueness is extra information and should not
be hidden inside a symbol.  Thus expressions such as \(H(t)\) or \(d_A(z)\)
are read as graph/predicate notation unless the text explicitly states that a
singleton branch has been selected.

\section*{Equation Audit}

\begingroup
\scriptsize
\setlength{\tabcolsep}{2pt}
\begin{longtable}{p{0.08\linewidth}p{0.06\linewidth}p{0.16\linewidth}p{0.31\linewidth}p{0.13\linewidth}p{0.12\linewidth}}
\toprule
Mills equation & Page & Section & Extracted TeX & Notebook symbol & Status\\
\midrule
\endhead
32.126 & 1536 & Hubble law &
\(\displaystyle v=HR\) &
redshift path premise & supporting\\
\addlinespace
32.36 & 1519 & Gravity / particle production &
\(\displaystyle r_g=\frac{2Gm_0}{c^2}\) &
\texttt{2 G m/c\^{}2} & matched\\
\addlinespace
32.38 & 1520 & Gravity metric &
\(\displaystyle
\dd\tau^2=-\left(1-\frac{2Gm_0}{c^2r}\right)\dd t^2
-\left(1-\frac{2Gm_0}{c^2r}\right)^{-1}\frac{\dd r^2}{c^2}
+\frac{r^2}{c^2}\dd\theta^2
+\frac{r^2\sin^2\theta}{c^2}\dd\phi^2
\) &
implicit & supporting\\
\addlinespace
32.43 & 1523 & Particle production &
\(\displaystyle
\tau=\frac{GM}{r^*c^2}=\frac{r_g}{c}=\frac{v_g}{c}t_i=t_i\Delta
\) &
time metric & supporting, OCR risk\\
\addlinespace
32.140 & 1542 & Expanding universe &
\(\displaystyle Q=\frac{c^3}{4\pi G}\) &
\texttt{c\^{}3/(4 Pi G)} & matched\\
\addlinespace
32.146 & 1544 & Minimum radius &
\(\displaystyle
\Re_{\min}=\left(\frac{\cdots}{4\pi\sigma eT^4}\right)^{1/2}
\) with printed result \(8.52\times 10^{27}\,\mathrm{m}\) &
\texttt{alephMin} & ambiguous norm.\\
\addlinespace
32.149 & 1545 & Oscillation period &
\(\displaystyle T_U=\frac{4\pi Gm_U}{c^3}\), but the printed numerical value is
\(3.10\times10^{19}\,\mathrm{s}\) for \(m_U\simeq2\times10^{54}\,\mathrm{kg}\) &
period in \texttt{aleph[t\_]} & mismatch: numeric convention uses \(2\pi Gm_U/c^3\)\\
\addlinespace
32.153 & 1546 & Radius evolution &
\(\displaystyle
X(t)=2.16\times10^{28}-1.86\times10^{28}
\cos\!\left(\frac{2\pi t}{9.83\times10^{11}\,\mathrm{yr}}\right)\,\mathrm{m}
\) &
\texttt{aleph[t\_]} & matched numeric\\
\addlinespace
32.154 & 1546 & Expansion rate &
\(\displaystyle \dot X=4\pi c\,
\sin\!\left(\frac{2\pi t}{9.83\times10^{11}\,\mathrm{yr}}\right)\) &
\texttt{alephRate[t\_]} & matched\\
\addlinespace
32.156 & 1547 & Hubble constant &
\(\displaystyle H(t)=\frac{\dot S}{ct}
=\frac{4\pi c\sin(\cdots)}{ct}\) &
\texttt{H[t\_]} & matched; denominator under audit\\
\addlinespace
32.158 & 1548 & Matter inventory &
\(\displaystyle
m_U(t)=\frac{m_U}{2}
\left(1+\cos\!\left(\frac{2\pi t}{2\pi Gm_U/c^3}\right)\right)
\) &
\texttt{mUBook[t\_]} & matched; commentary conflict\\
\addlinespace
32.165 & 1551 & Temp./redshift &
\(\displaystyle
\lambda_\infty=\lambda(r)\left(1+\frac{2Gm_U}{rc^2}\right)
\) &
\texttt{zGUTCP} & endpoint correction\\
\addlinespace
p.1579 & 1579 & CMBR structure scale &
\(\displaystyle L_\ell=\frac{2}{\ell}\,2\pi r_{\mathrm{sphere}}
=\frac{4\pi r_{\mathrm{sphere}}}{\ell}\) &
\texttt{cmbrScale} & stated; not BAO\\
\bottomrule
\end{longtable}
\endgroup

\section*{GUTCP \(z\): Denominator Audit}

[FORCED] The observational redshift is defined by the wavelength ratio
\[
1+z=\frac{\lambda_o}{\lambda_e}.
\]
Mills gives a Hubble-law structure in Eq. (32.126), \(v=HR\), and then gives
the oscillatory radius and rate in Eqs. (32.153)--(32.154).  Page 1540--1544
frames the mechanism as mass-energy conversion causing spacetime expansion:
matter creation contracts spacetime, matter-to-energy release expands it, and
redshift/blueshift time stamps the extent of spacetime expansion between
observer and observed.  Therefore the denominator in \(H=\dot R/R\) must be
the spacetime scale or extent being compared.

Using the notebook and the numerical convention of Eqs. (32.153)--(32.156),
\begin{align}
\theta(t) &= \omega t,&
\omega &= \frac{2\pi}{T_U},&
T_U &= \frac{2\pi GM_0}{c^3},\\
\RG(t) &= \frac{2GM_0}{c^2}
+\frac{4\pi GM_0}{c^2}\left[1-\cos\theta(t)\right],&
\dot{\RG}(t)&=4\pi c\,\sin\theta(t).
\end{align}
The printed Eq. (32.149) instead gives \(T_U=4\pi GM_0/c^3\); that
factor-of-two choice remains a convention switch.

[CHOSEN] Mills's Eq. (32.156) replaces the Hubble-law denominator by \(ct\):
\begin{equation}
H_{ct}(t)=\frac{\dot{\RG}(t)}{ct}
=\frac{4\pi\sin\theta(t)}{t}.
\label{eq:h-ct}
\end{equation}
This is the step that creates a damped sinc-like graph from an oscillatory
radius.  Its associated path-redshift relation is
\begin{equation}
\boxed{
\ZRel^{ct}(t_e,t_o,z)
\Longleftrightarrow
1+z
=
\exp\!\left[
\int_{t_e}^{t_o} H_{ct}(t)\,\dd t
\right]
}
\label{eq:zgutcp-ct}
\end{equation}
or, for the numerical period convention,
\begin{equation}
\boxed{
\ZRel^{ct}(t_e,t_o,z)
\Longleftrightarrow
1+z
=
\exp\!\left[
4\pi\left(\operatorname{Si}(\omega t_o)
-\operatorname{Si}(\omega t_e)\right)
\right],
}
\label{eq:zgutcp-ct-si}
\end{equation}
where \(\operatorname{Si}(x)=\int_0^x(\sin u/u)\,\dd u\).  This preserves the
notebook calculation, but it is not forced by the arrow-of-time construction.

[FORCED] If Eq. (32.126) is applied to the whole oscillatory GUTCP radius, then
\[
H_{\RG}(t)=\frac{\dot{\RG}(t)}{\RG(t)},\qquad
\dd\ln(1+\zG)=H_{\RG}(t)\,\dd t.
\]
The integral is exact:
\begin{equation}
\boxed{
\ZRel^{\RG}(t_e,t_o,z)
\Longleftrightarrow
1+z=\frac{\RG(t_o)}{\RG(t_e)}.
}
\label{eq:zgutcp-radius}
\end{equation}
This expression has the same algebraic shape as a FLRW scale-factor ratio, but
here it is derived from Mills's \(v=HR\) and his \(\RG(t)\), not assumed from
FLRW geometry.

[CHOSEN] The stronger reading of page 1542 is that redshift time stamps the
\emph{extent of spacetime expansion} from the minimum-radius bounce, not the
entire radius including the finite gravitational-radius offset.  Define
\begin{equation}
\EG(t)=\RG(t)-\RG_{\min},\qquad
\RG_{\min}=\frac{2GM_0}{c^2}.
\end{equation}
Then
\[
H_{\EG}(t)=\frac{\dot{\EG}(t)}{\EG(t)}
=\frac{\dot{\RG}(t)}{\RG(t)-\RG_{\min}},
\]
and the corrected expansion-extent redshift is
\begin{equation}
\boxed{
\ZRel^{\EG}(t_e,t_o,z)
\Longleftrightarrow
1+z=
\frac{\RG(t_o)-\RG_{\min}}{\RG(t_e)-\RG_{\min}}.
}
\label{eq:zgutcp-extent}
\end{equation}
This removes the artificial \(1/t\) damping introduced by \(ct\), and it keeps
Mills's claim that redshift/blueshift records the extent of expansion.  It is
also the branch form implemented as the Python default
\texttt{expansion\_extent}.

\section*{Endpoint Clock Correction}

The second redshift-relevant ingredient is Mills's explicit gravitational
wavelength statement, Eq. (32.165):
\begin{equation}
\lambda_\infty=\lambda(r)B(r),
\qquad
B(r)=1+\frac{2Gm_U}{rc^2}.
\label{eq:mills-redshift-factor}
\end{equation}

[FORCED] If Eq.~\eqref{eq:mills-redshift-factor} is accepted as a wavelength
shift law and if two finite \(r\)-spheres are compared through the same
\(\lambda_\infty\) reference, then the algebraic endpoint clock ratio is
\begin{equation}
\lambda_o B(r_o)=\lambda_e B(r_e),
\qquad
\ZRel^B(t_e,t_o,z)\Longleftrightarrow
1+z=\frac{\lambda_o}{\lambda_e}
=\frac{B(r_e)}{B(r_o)}.
\end{equation}
Removing Eq. (32.165) removes the wavelength-shift law; removing the common
\(\lambda_\infty\) comparison removes the finite-radius ratio.

With \(r=\RG(t)\), the endpoint factor is
\begin{equation}
\boxed{
\ZRel^{B}(t_e,t_o,z;M_z)
\Longleftrightarrow
1+z
=
\frac{
1+\dfrac{2G\,M_z(t_e)}{c^2\RG(t_e)}
}{
1+\dfrac{2G\,M_z(t_o)}{c^2\RG(t_o)}
}
}
\label{eq:zgutcp-endpoint}
\end{equation}
where \(M_z(t)\) is the mass parameter used in the redshift factor.  The most
literal Eq. (32.165) choice is \(M_z(t)=M_0\).  The Eq. (32.158) choice is
\begin{equation}
M_z(t)=m_U(t)=\frac{M_0}{2}\left(1+\cos\theta(t)\right).
\end{equation}

[CHOSEN] If the Eq. (32.165) endpoint clock factor is treated as separate from
the expansion-extent stretching, the combined redshift relation is
\begin{equation}
\boxed{
\ZRel^{\EG B}(t_e,t_o,z;M_z)
\Longleftrightarrow
1+z
=
\frac{\RG(t_o)-\RG_{\min}}{\RG(t_e)-\RG_{\min}}
\frac{
1+\dfrac{2G\,M_z(t_e)}{c^2\RG(t_e)}
}{
1+\dfrac{2G\,M_z(t_o)}{c^2\RG(t_o)}
}.
}
\label{eq:zgutcp-combined}
\end{equation}
The Python artifact exposes the denominator choices explicitly:
\texttt{expansion\_extent}, \texttt{radius\_scale}, \texttt{mills\_ct\_path},
\texttt{mills\_endpoint}, and combined endpoint modes.  The
\texttt{expansion\_extent} mode is the selected GUTCP redshift relation
counterpart to the FLRW \(z\) parameter; the endpoint factor is a printed Mills
correction relation that can be included or audited separately.

\section*{Branch Inversion}

[FORCED] For fixed \(t_o\), a redshift law is a relation on
\((t_e,z)\), not automatically a singleton graph in either direction.  For a selected
redshift relation \(\ZRel^\star\), define the emission-time fiber
\begin{equation}
\mathcal{T}_{z}^{\star}(z;t_o)=
\left\{t_e<t_o:\ \ZRel^\star(t_e,t_o,z)\right\}.
\label{eq:emission-time-fiber}
\end{equation}
This set may contain zero, one, or many emission times.  A branch condition
\(\mathfrak{B}_{\mathrm{GUTCP},b}(t_e,t_o,z)\) restricts the fiber to
\begin{equation}
\mathcal{T}_{z,b}^{\star}(z;t_o)=
\left\{t_e<t_o:\ \ZRel^\star(t_e,t_o,z)\wedge
\mathfrak{B}_{\mathrm{GUTCP},b}(t_e,t_o,z)\right\}.
\label{eq:branch-restricted-fiber}
\end{equation}
Only when this restricted fiber is a singleton is it harmless to write
\(t_b(z)\) as shorthand for its sole member.

For a radial light-front relation, the corresponding light-travel-distance
relation is
\begin{equation}
\DRel^\star(z,D;t_o)
\Longleftrightarrow
\exists t_e\,
\left[
t_e\in\mathcal{T}_{z}^{\star}(z;t_o)
\wedge D=c(t_o-t_e)
\right].
\label{eq:light-distance-relation}
\end{equation}
This relation is the necessary replacement for a single FLRW \(D(z)\).
Oscillation does not remove distance-redshift solutions; it changes the
cardinality of the fibers and therefore makes branch information explicit.

\section*{Blueshift Threshold Relations}

[FORCED] If \(t_o\) is measured from the start of the current expansion and a
radial light-front relation is used,
\begin{equation}
\mathfrak{D}_{\mathrm{light}}(t_e,t_o,D)
\Longleftrightarrow
D=c(t_o-t_e),
\label{eq:radial-light-distance}
\end{equation}
then the threshold at which the emission event leaves the current-expansion
branch is
\begin{equation}
\mathfrak{D}_{\mathrm{enter-blue}}(D_\star,t_o)
\Longleftrightarrow
D_\star=c\,t_o.
\label{eq:blue-segment-entry-threshold}
\end{equation}
For \(D<c\,t_o\), the emission time satisfies \(0<t_e<t_o\), so the path lies
inside the current expansion branch.  For \(D>c\,t_o\), the relation gives
\(t_e<0\), so a pre-current-expansion segment has entered the path.  That is
the first threshold at which a contraction/blueshift segment can become part
of the redshift relation.

[CHOSEN] This threshold is not the same as the threshold for \emph{net}
endpoint blueshift.  For the even oscillatory endpoint relations
\(\RG(-t)=\RG(t)\) and \(\EG(-t)=\EG(t)\), the endpoint-only neutral condition
is \(t_e=-t_o\), hence
\begin{equation}
\mathfrak{D}_{\mathrm{net-neutral}}(D_\star,t_o)
\Longleftrightarrow
D_\star=2c\,t_o.
\label{eq:net-blue-threshold}
\end{equation}
Under that endpoint-only symmetry, \(D>2c\,t_o\) gives an emission event whose
oscillatory radius/extent exceeds the observation value, so the endpoint ratio
is net blueshift.  Between \(c\,t_o\) and \(2c\,t_o\), a blueshift segment may
enter the path while the endpoint ratio can still be net redshift.

[FORCED] These thresholds are branch-admissibility relations.  If an
independent observational relation supplies or bounds \(D\), then
\[
D<c\,t_o
\]
selects the current-expansion branch and excludes branches requiring
\(t_e<0\).  Conversely,
\[
D>c\,t_o
\]
excludes a current-expansion-only branch and forces a segmented relation that
admits a pre-current-expansion portion of the path.

[CHOSEN] If the symmetric endpoint-only rule is the branch rule being tested,
the observed sign of \(z\) adds a second pruning condition:
\begin{equation}
\mathfrak{B}_{\mathrm{sign}}(D,t_o,s)
\Longleftrightarrow
\begin{cases}
D<2c\,t_o, & s=\mathrm{red},\\
D=2c\,t_o, & s=\mathrm{neutral},\\
D>2c\,t_o, & s=\mathrm{blue}.
\end{cases}
\label{eq:blue-sign-branch-admissibility}
\end{equation}
Thus a positive observed redshift can rule out a candidate branch whose
endpoint-symmetric relation already lies beyond the net-blue threshold.  This
does not make the full distance-redshift relation unique by itself; it removes
inadmissible members from the fiber before any singleton branch is selected.

[WEAKEST LINK] The expansion-extent relation has \(\EG(0)=0\).  A light path
that crosses the exact minimum-radius boundary therefore requires an additional
bounce or matching relation before the segmented redshift product is physically
well defined.  The whole-radius relation avoids this zero but includes the
finite gravitational-radius offset.

\section*{Koksbang--Heinesen FLRW-Consistency Standard}

The local reference paper by Koksbang and Heinesen, stored under
\texttt{docs/cites/} as the model-independent generalized-FLRW consistency PDF,
sets a useful external standard for a fair comparison.  Its model-independent
pipeline does not begin with a cosmological parameter fit.  It reconstructs
redshift-indexed observable relations, conventionally written in graph
shorthand as
\begin{equation}
d_A(z),\qquad d_A'(z),\qquad d_A''(z),\qquad
H_{\parallel}(z),\qquad H_{\parallel}'(z),
\label{eq:kh-input-bundle}
\end{equation}
where \(d_A\) is the angular diameter distance and \(H_{\parallel}\) is the
line-of-sight expansion rate
\begin{equation}
H_{\parallel}=\frac{1}{3}\theta-e^\mu a_\mu+e^\mu e^\nu\sigma_{\mu\nu}.
\label{eq:kh-line-of-sight-h}
\end{equation}
The paper then evaluates generalized FLRW consistency diagnostics over the
overlap of the supernova and BAO reconstructions, with the currently useful
domain beginning at \(z=0.38\) and extending to order \(z\simeq2\).

[FORCED] BAO is not a direct model-independent \(d_A\) datum.  The measured
object is a clustering-feature relation in angle/redshift space; the usual
compressed reductions are ratios or products such as \(D_M(z)/r_d\),
\(H_{\parallel}(z)r_d\), or \(D_V(z)/r_d\).  Transverse BAO therefore becomes
an angular-ruler constraint only after a sound-horizon/ruler relation and a
distance dictionary are admitted.  In a GUTCP comparison it belongs on the
conditional data-model side, not as the missing GUTCP angular-diameter
propagation relation itself.

[FORCED] Mills's page-1579 treatment is a different claimed ruler relation,
not a BAO reduction.  He takes the present Earth \(r\)-sphere radius as
\[
r_{\mathrm{sphere}}=14.02\times10^9\ \mathrm{ly}
\]
and the corresponding angular view as
\[
A_{\mathrm{view}}=2\pi r_{\mathrm{sphere}}
=88.09\times10^9\ \mathrm{ly}.
\]
For a CMBR multipole \(\ell\), the stated sky fraction is \(2/\ell\), so the
large-structure scale relation is
\begin{equation}
\mathfrak{M}_{\mathrm{CMBR}}
(\ell,r_{\mathrm{sphere}},L)
\Longleftrightarrow
L=\frac{2}{\ell}A_{\mathrm{view}}
=\frac{4\pi r_{\mathrm{sphere}}}{\ell}.
\label{eq:mills-cmbr-structure-scale}
\end{equation}
This gives \(L_{15}=11.7\times10^9\) ly, \(L_{135}=1.3\times10^9\) ly, and
\(L_{700}=0.25\times10^9\) ly, matching the examples Mills assigns to the
Hercules--Corona Borealis Great Wall, the Big Ring, and smaller higher-order
structures.  The relation is therefore a GUTCP/CMBR multipole structure-scale
claim.  It should not be collapsed into the BAO transverse-scale data product.
Testing it against observed arcs or rings still requires the relevant object
selection, angular extent, redshift extent, and observational distance
dictionary to be stated as relations.

In relation notation, the angular-diameter observation is at least
\begin{equation}
\mathfrak{A}_{\mathrm{obs}}
\left(z,d_A,\dd A_{\mathrm{emit}},\dd\Omega_{\mathrm{obs}}\right)
\Longleftrightarrow
d_A^2=\frac{\dd A_{\mathrm{emit}}}{\dd\Omega_{\mathrm{obs}}}.
\label{eq:angular-diameter-observation}
\end{equation}
This relation can be stated without FLRW.  What is still needed for GUTCP is
the constitutive relation that connects
\((\dd A_{\mathrm{emit}},\dd\Omega_{\mathrm{obs}})\) to a selected GUTCP light
path, phase, mass convention, and branch.

Two of the diagnostic targets are especially useful for auditing a non-FLRW
proposal.  The integrated CBL relation is written in singleton-graph shorthand as
\begin{equation}
O(z)=H_{\parallel}(z)^2D'(z)^2-1,
\qquad
O \xrightarrow{\mathrm{FLRW}} \Omega_{k,0}H_0^2,
\label{eq:kh-o}
\end{equation}
where \(D\) is the distance variable used in the CBL integral relation.  The
density/\(\Lambda\)CDM diagnostic is
\begin{equation}
M(z)=
\frac{2H_{\parallel}(z)^2}{3d_A(z)(1+z)}
\left[
-d_A'(z)
+\left(\frac{2}{1+z}+\frac{H_{\parallel}'(z)}{H_{\parallel}(z)}\right)d_A(z)
+d_A''(z)
\right],
\label{eq:kh-m}
\end{equation}
with \(M\to\Omega_{m,0}H_0^2\) in \(\Lambda\)CDM under the paper's stated
assumptions.  The full generalized CBL statistic \(C(z)\) similarly combines
\(d_A\), \(H_{\parallel}\), their derivatives, optical shear/expansion, and
Ricci focusing; the FLRW null expectation is \(C=0\).

[FORCED] Applying that standard to GUTCP requires more than a redshift
relation.  A branch relation must be selected so that
\[
t_e\in\mathcal{T}_{z,b}^{\star}(z;t_o),
\]
and then the branch-restricted GUTCP quantities must be joined with
observable relations.  For example,
\begin{align}
\mathfrak{H}_{\parallel,\mathrm{GUTCP},b}^{\star}(z,H;t_o)
&\Longleftrightarrow
\exists t_e\,
\left[
t_e\in\mathcal{T}_{z,b}^{\star}(z;t_o)
\wedge H=H_{\parallel}(t_e)
\right],
\\
\mathfrak{A}_{\mathrm{GUTCP},b}^{\star}(z,d_A;t_o)
&\Longleftrightarrow
\exists t_e\,
\left[
t_e\in\mathcal{T}_{z,b}^{\star}(z;t_o)
\wedge
\mathfrak{A}_{\mathrm{obs}}(\cdots)
\wedge
\mathfrak{L}_{\mathrm{GUTCP},b}(\cdots)
\right].
\end{align}
Here \(\mathfrak{L}_{\mathrm{GUTCP},b}\) denotes the missing GUTCP
light-bundle/angular propagation relation.  The corrected expansion-extent
\(\ZRel^{\EG}\) constrains the redshift-time relation, but it does not by
itself constrain \(\mathfrak{L}_{\mathrm{GUTCP},b}\).

[FORCED] The frontier relation is open, but it is not unconstrained.  Pages
1569, 1577--1579, and the page-1602 footnote give information that any candidate
\(\mathfrak{L}_{\mathrm{GUTCP},b}\) must either incorporate or explicitly
exclude:
\[
\Delta T_E(\ell)=
C_{\mathrm{effThompson}}\,77\,
\mathrm{sinc}\!\left(\frac{\pi}{140}(\ell+197)\right)\,\mu\mathrm{K},
\]
\[
\Delta T_B(\ell)=
r^{1/2}C_{\mathrm{effThompson}}\,77\,
\mathrm{sinc}\!\left(\frac{\pi}{140}(\ell+197+70)\right)\,\mu\mathrm{K},
\qquad
r^{1/2}=\frac{\Delta\aleph}{ct}.
\]
Mills identifies these as polarization/lensing constraints: scattering
produces E-mode correlation multipoles, gravitational lensing and accelerating
spacetime convert E-mode structure into B-mode structure, and the large-scale
matter distribution inherits CMBR multipole/lensing patterns.  Together with
Eq.~\eqref{eq:mills-cmbr-structure-scale}, these pages constrain the possible
angular/light-bundle closure but still do not supply the map
\[
(t_e,t_o,\mathrm{branch},\mathrm{source\ geometry})
\longmapsto
(\dd A_{\mathrm{emit}},\dd\Omega_{\mathrm{obs}}).
\]
Page 1569 adds a separate constraint and a separate risk: Mills uses CMBR
structures of about \(1^\circ\) as evidence for nearly flat geometry since the
commencement of expansion, but the same sentence ties the geometry claim to a
commencement of expansion ``about 10 billion years ago.''  Since the notebook
calibration cannot generally accept \(t_{\mathrm{now}}=10^{10}\) yr as an
input without overconstraining the \(H_0\)/CMB temperature solve, any
\(\mathfrak{L}_{\mathrm{GUTCP},b}\) derived from that geometry discussion must
separate the near-flat/angular-structure claim from the fixed-10-Gyr
approximation.

The page-1602 footnote adds an observational-reference constraint.  Mills
identifies a practical reference for absolute space at rest as a bright
celestial body with zero translational velocity within its \(r\)-sphere, or
with that velocity component corrected.  He then proposes a Cepheid of the
corresponding calculated age/distance relative to the current \(r\)-sphere as
an approximate point at rest on a selected \(r\)-sphere, and states that
changes in angular diameter over the pulsation cycle combined with
spectroscopic radial velocity determine distance quasi-geometrically.  A
candidate \(\mathfrak{L}_{\mathrm{GUTCP},b}\) therefore cannot use an arbitrary
observer frame: it must either encode this absolute-rest/Cepheid calibration
relation or explicitly explain why it is not part of the optical dictionary.

\begingroup
\scriptsize
\setlength{\tabcolsep}{3pt}
\begin{longtable}{p{0.12\linewidth}p{0.32\linewidth}p{0.32\linewidth}p{0.14\linewidth}}
\toprule
Source & Relation fragment & Constraint on \(\mathfrak{L}_{\mathrm{GUTCP}}\) & Status\\
\midrule
\endhead
p.1569 & BOOMERANG \(1^\circ\) CMBR structures and nearly flat geometry since
commencement of expansion & Candidate angular closure must decide whether this
is a curvature/flatness constraint, while not importing fixed \(10^{10}\) yr as
a calibration input & constraining input with calibration hazard\\
\addlinespace
p.1602 footnote & Absolute-rest reference object and Cepheid angular-diameter
plus spectroscopic-radial-velocity distance calibration & Candidate optical
dictionary must specify the \(r\)-sphere rest-frame/velocity-correction rule
and how quasi-geometrical Cepheid distances anchor the scale & constraining
input\\
\addlinespace
p.1577, Eq. 32.203 & E-mode CMBR multipole pattern from Thompson scattering &
Any optical relation including polarization must preserve the stated
\(\ell\)-indexed E-mode phase and amplitude convention & constraining input\\
\addlinespace
p.1577, Eqs. 32.204--32.205 & B-mode pattern shifted by \(70\) in \(\ell\)
with amplitude ratio \(\Delta\aleph/(ct)\) & Candidate lensing/accelerating
spacetime closure must explain the E-to-B conversion and amplitude ratio &
constraining input\\
\addlinespace
p.1578 & Gravitational lensing produces rings/arcs and is linked to matter
formation patterns & Candidate light-bundle relation must decide how lensing
maps photon bundles into observed angular structures & constraining input\\
\addlinespace
p.1579 & CMBR multipole scale \(L_\ell=4\pi r_{\mathrm{sphere}}/\ell\) &
Large-ring/arc comparisons should use the CMBR multipole ruler relation, not
BAO transverse scale by default & stated relation\\
\bottomrule
\end{longtable}
\endgroup

\begingroup
\scriptsize
\setlength{\tabcolsep}{2pt}
\begin{longtable}{p{0.14\linewidth}p{0.28\linewidth}p{0.28\linewidth}p{0.20\linewidth}}
\toprule
Required relation & Paper standard & Current GUTCP artifact & Fair evaluation\\
\midrule
\endhead
\(\ZRel\) & Observed redshift relation for \(d_A\) and \(H_{\parallel}\) &
\(\ZRel(t_e,t_o,z)\) stated after choosing denominator/endpoint convention &
conditional, branch dependent\\
\addlinespace
\(\mathcal{T}_{z,b}\) & Restrict redshift fibers to index model quantities &
Implemented as multi-valued emission-time fibers & requires branch/window bits\\
\addlinespace
\(\HRel\) & Line-of-sight expansion rate relation, not merely a named
``Hubble constant'' & \(H_{\EG}\), \(H_{\RG}\), or Mills \(H_{ct}\) joined to a
selected branch fiber & partial; convention must be declared\\
\addlinespace
\(\partial_z\HRel\) & Redshift derivative relation for \(C\) and \(M\) &
reducible only after a stable branch relation & not yet paper-equivalent\\
\addlinespace
\(\ARel\) & Angular-diameter observational relation reconstructed from
supernovae & general relation stated in Eq.~\eqref{eq:angular-diameter-observation};
GUTCP angular-propagation relation not stated & relation gap\\
\addlinespace
\(\partial_z\mathfrak{A},\partial_z^2\mathfrak{A}\) & Required distance-derivative relations
& not reducible until \(\ARel\) is joined to a GUTCP branch/angular relation &
relation gap\\
\addlinespace
\(\mathrm{BAO}\) ruler & BAO constrains \(D_M/r_d\),
\(H_{\parallel}r_d\), or \(D_V/r_d\), not bare \(d_A\) & can be joined only as
a conditional ruler/data relation & not a model-independent GUTCP \(d_A\)
datum\\
\addlinespace
\(C,O,M\) & Diagnostic FLRW/general-spacetime consistency relations with
uncertainty bands & cannot be fairly reduced from \(H(t)\) and light-travel
distance alone & not yet evaluable at full standard\\
\bottomrule
\end{longtable}
\endgroup

[CHOSEN] The Python artifact therefore exposes a readiness table and
branch-specific surrogate observable relations rather than forcing a numerical
Koksbang--Heinesen \(C/O/M\) reduction.  For example, from a selected branch
one may state
\[
\mathfrak{D}_{\mathrm{GUTCP},\mathrm{light},b}^{\star}(z,D;t_o)
\Longleftrightarrow
\exists t_e\,
\left[
t_e\in\mathcal{T}_{z,b}^{\star}(z;t_o)
\wedge D=c(t_o-t_e)
\right]
\]
and a candidate
\(\mathfrak{H}_{\parallel,\mathrm{GUTCP},b}^{\star}(z,H;t_o)\).  But
\[
\mathfrak{D}_{\mathrm{GUTCP},\mathrm{light},b}^{\star}
\ne
\mathfrak{A}_{\mathrm{GUTCP},b}^{\star}
\]
unless an additional GUTCP angular-size propagation law is supplied.  Treating
light-travel distance as angular diameter distance would import an
observational dictionary not present in the audited material.

[FORCED] A fair comparison to the paper is therefore a readiness verdict, not a
premature rejection.  GUTCP now has candidate redshift relations and
branch-restricted fibers.  It does not yet meet the paper's full standard
because the GUTCP angular-diameter relation, its derivative relations, the
optical shear/focusing relations, and an uncertainty relation/pipeline are not
yet joined to the current artifact.

[CHOSEN] In MDL terms, the information needed to select a GUTCP branch/window
is part of the model description length.  That cost should be compared against
the information needed by FLRW models to describe the observed departures in
angular-size and distance diagnostics.  The paper's standard is suitable for
that comparison because it asks for observable relations and residuals, not for
membership in the FLRW class.

\section*{Appendix A: \(H_0\) Target Sweep}

The notebook line
\[
\texttt{H0GUTCP = H[tNow]}
\]
defines a model quantity:
\begin{equation}
H_{0,\mathrm{GUTCP}}(t_{\mathrm{now}},T_{\mathrm{start}})
=H(t_{\mathrm{now}};M_0(T_{\mathrm{start}})).
\end{equation}
It is not, by itself, an independent prediction of the observed Hubble
constant.  The later Mathematica equality
\[
\texttt{(H0GUTCP/.quantitiesSubs) == QuantityMagnitude[UnitConvert[Subscript[H,0]]]}
\]
is a calibration condition.  To make this explicit, replace the Wolfram entity
lookup with a named target:
\begin{align}
H_{0,\mathrm{GUTCP}}(t_{\mathrm{now}},T_{\mathrm{start}})
&=H_{0,\mathrm{target}},\\
T_U(t_{\mathrm{now}};T_{\mathrm{start}})
&=T_{\mathrm{CMB,obs}}.
\end{align}
This is why the original notebook solved for \(t_{\mathrm{now}}\) instead of
using Mills's approximate \(10^{10}\) yr value as an input.  If
\(t_{\mathrm{now}}\) is fixed, then \(T_{\mathrm{start}}\) is the only remaining
unknown, but the notebook has two empirical equations to satisfy:
\[
H_{0,\mathrm{GUTCP}}(T_{\mathrm{start}})=H_{0,\mathrm{target}},
\qquad
T_U(T_{\mathrm{start}})=T_{\mathrm{CMB,obs}}.
\]
Those two constraints are generally inconsistent for one unknown.  Leaving
\((t_{\mathrm{now}},T_{\mathrm{start}})\) free gives a determined two-equation,
two-unknown calibration problem rather than an overconstrained one.

The sweep variable is therefore \(H_{0,\mathrm{target}}\), commonly entered as
\texttt{H0TargetKmSecMpc}.  Each selected target gives a calibrated pair
\((t_{\mathrm{now}},T_{\mathrm{start}})\).  This makes the Hubble-tension scan a
model-fitting exercise, not a claim that the GUTCP equations alone selected the
empirical \(H_0\).

\section*{Appendix B: Observational Relation Gap}

The current notebook states time-parametrized relations:
\[
H(t),\qquad \mathrm{LightRadius}(t),\qquad T_U(t),\qquad m_U(t),\qquad P_U(t).
\]
It does not yet state a complete observational relation
\[
t \longleftrightarrow z_{\mathrm{GUTCP}}
\]
that is independent of FLRW and reduced by branch constraints.  In particular,
one should not assume the whole-radius relation
\[
\ZRel^{\RG}(t_e,t_o,z)\Longleftrightarrow
1+z=\frac{\RG(t_o)}{\RG(t_e)}
\]
merely by analogy with FLRW.  The denominator audit above derives two GUTCP
relations from Mills's own variables: the whole-radius relation
Eq.~\eqref{eq:zgutcp-radius}, and the expansion-extent relation
Eq.~\eqref{eq:zgutcp-extent}.  The latter is selected here because page 1542
describes redshift/blueshift as time stamps of the \emph{extent of spacetime
expansion} between observer and observed.

Eq.~\eqref{eq:zgutcp-extent} states a GUTCP-native redshift relation for a
specified emission time, observation time, and redshift value.  Oscillation does
not imply that the distance-redshift relation has no solution; it implies that
the relevant fibers may have multiple members.  A complete observational
relation must include the relevant branch, observation window, and light path
through phases where ancient light may undergo partial blueshifting as well as
redshifting.  Only after such constraints are joined can external diagnostics be
stated as relations:
\begin{align}
\mathfrak{H}_{\parallel,\mathrm{GUTCP},b}(z,H;t_o)
&\Longleftrightarrow
\exists t_e\,
\left[
t_e\in\mathcal{T}_{z,b}^{\star}(z;t_o)\wedge H=H_{\parallel}(t_e)
\right],
\\
\mathfrak{D}_{\mathrm{GUTCP},b}(z,D;t_o)
&\Longleftrightarrow
\exists t_e\,
\left[
t_e\in\mathcal{T}_{z,b}^{\star}(z;t_o)\wedge D=c(t_o-t_e)
\right].
\end{align}
This is the point at which comparisons to non-FLRW observational frameworks,
including Heinesen--Clifton style line-of-sight expansion and distance
diagnostics, become well-posed.  If \(\mathcal{T}_{z,b}^{\star}\) is later
shown to be a singleton fiber over a chosen domain, then the familiar
singleton-graph notation is merely a compression of these relations.

\section*{Audit Conclusion}

[FORCED] Mills's page 1540--1544 arrow-of-time argument makes spacetime
expansion the dynamic variable.  Eq. (32.126) then requires the denominator of
\(H=\dot R/R\) to be the compared spacetime scale or extent.  The later
substitution \(R=ct\) in Eq. (32.156) is not forced by that derivation and
produces a damped sinc-like \(H_{ct}(t)\) from an otherwise oscillatory
\(\RG(t)\).

[CHOSEN] The corrected GUTCP counterpart to FLRW \(z\) selected here is the
expansion-extent relation
\[
\ZRel^{\EG}(t_e,t_o,z)
\Longleftrightarrow
1+z=
\frac{\RG(t_o)-\RG_{\min}}{\RG(t_e)-\RG_{\min}},
\]
because Mills describes redshift/blueshift as a timestamp of the extent of
spacetime expansion.  The whole-radius relation is also exposed as a
convention, but it includes the finite gravitational-radius offset rather than
just the expansion extent.

[FORCED] Eq.~\eqref{eq:zgutcp-endpoint} is separately forced as a finite
endpoint clock ratio after accepting Mills's Eq. (32.165), a common
\(\lambda_\infty\) reference, and the substitution \(r=\RG(t)\).

[CHOSEN] The Python artifact defaults to the notebook/numerical period convention
\(T_U=2\pi GM_0/c^3\), because that is the convention that reproduces Eq. (32.156)
and the quoted \(H_0\simeq78.5\,\mathrm{km\,s^{-1}\,Mpc^{-1}}\) at
\(t=10^{10}\) yr when used as Mills's worked example.  The calibrated notebook
does not use \(10^{10}\) yr as an input because doing so would overconstrain the
two empirical conditions.  That calibration remains a \(\texttt{mills\_ct}\)
calibration unless the denominator convention is changed and the fit is repeated.
The Python artifact exposes the printed Eq. (32.149) convention as an option.

[FORCED] Under the Koksbang--Heinesen consistency standard, a fair external
evaluation must use angular-diameter, distance-derivative, line-of-sight
expansion, and expansion-derivative relations.  The current GUTCP artifact
states branch-specific redshift fibers, candidate line-of-sight expansion
relations, and light-travel-distance relations.  It has not yet joined these
to the GUTCP angular-diameter relation required by the paper, because the
needed light-propagation and angular-size relation has not been stated.

[WEAKEST LINK] The first weak link is Mills's unforced \(ct\) denominator in
Eq. (32.156).  The second is branch selection: Mills does not provide enough
material in the audited 2023 pages to select a singleton inverse fiber
\(\mathcal{T}_{z,b}^{\star}\) for light that crosses the oscillatory history.
The third is the angular-diameter relation.  The corrected redshift relation
can be stated, but a paper-equivalent \(C/O/M\) comparison remains conditional
until the branch relation, line-of-sight expansion relation, angular-diameter
relation, and uncertainty relation/pipeline are joined.

\end{document}
